\topic{Abstract.} Can a network learn to organize its computation more efficiently than a human-designed architecture? We tested input-dependent depth, grouping, expert routing and memory in eight experiments trained from scratch. On Boolean programs, adaptive depth reached 85.53\% accuracy and random depth 95.66\% at nearly identical estimated inference FLOPs. The clearest benefit came from choosing which information to retain and group. On a four-class synthetic task, grouping plus memory reached 100\% across three seeds. The full controller reached 91.17\%, and a policy written with knowledge of the task rules reached 99.84\%. Shuffling learned actions reduced full-model accuracy by 50.75 percentage points. Independent experiments recovered latent symbol families and retained relevant records under distribution shift. Uniform random expert routing matched learned routing on the task mixture. On natural text a five-minute, 118K-parameter decoder beat a smoothed bigram and generated dialogue-like bytes, yet it exited at one pass on every evaluated window. The experiments support several useful task-specific learned decisions. A general efficiency advantage of learning the entire computational organization remains unestablished.
